\subsection{短期冲击模型（0--60 天）}

短期冲击模型对应赛题任务 1。模型以赛题给定的供应中断、SPR、商业库存、绕道运输和恐慌情绪为基础，并纳入能够通过消融实验和敏感性分析检验的价格形成因素。其目标是在附件冲突窗口内解释 0--60 天真实价格路径，并为后续中长期油价调节模型提供统一机制基础。

\begin{figure}[htbp]
\centering
\begin{tikzpicture}[
  x=1cm,
  y=1cm,
  layer/.style={draw=OilBorder, fill=PaperNoteBack, rounded corners=1pt, align=center, minimum height=0.78cm, text width=3.15cm, font=\small\bfseries},
  item/.style={draw=OilBorder, fill=white, rounded corners=1pt, align=center, minimum height=0.58cm, text width=2.72cm, font=\scriptsize},
  arrow/.style={-{Latex[length=2.1mm]}, line width=0.45pt, draw=OilGray},
  feedback/.style={-{Latex[length=2.0mm]}, line width=0.4pt, draw=OilGray, dashed}
]
\node[layer] (physical) at (0,0) {物理供需层};
\node[layer] (pricing) at (4.45,0) {市场价格形成层};
\node[layer] (path) at (8.90,0) {日度价格路径};
\node[item] (gap) at (0,-1.28) {供应中断\\需求收缩};
\node[item] (buffer) at (0,-2.12) {SPR\\库存与绕道};
\node[item] (risk) at (4.45,-1.28) {风险溢价\\制度不确定性};
\node[item] (relief) at (4.45,-2.12) {缓冲确认\\预期修复};
\node[item] (eval) at (8.90,-1.70) {拟合、基准\\消融检验};
\draw[arrow] (gap) -- (physical);
\draw[arrow] (buffer) -- (physical);
\draw[arrow] (physical) -- node[above, font=\scriptsize] {剩余缺口 \(G_t\)} (pricing);
\draw[arrow] (risk) -- (pricing);
\draw[arrow] (relief) -- (pricing);
\draw[arrow] (pricing) -- node[above, font=\scriptsize] {目标价格 \(P_t^\ast\)} (path);
\draw[arrow] (path) -- (eval);
\draw[feedback] (eval.south west) .. controls (6.6,-3.0) and (2.4,-3.0) .. node[below, font=\scriptsize] {反馈校验} (physical.south);
\end{tikzpicture}
\caption{短期冲击模型结构}
\end{figure}

\keypoint{主模型分为两层：物理供需层负责回答“缺口还剩多少”，市场价格形成层负责回答“市场如何给剩余缺口和制度风险定价”。两层分开写，使价格平台能够被分解为物理缓冲、风险溢价和预期修复的共同结果。}

\begin{figure}[htbp]
\centering
\begin{tikzpicture}[
  node distance=1.05cm and 1.05cm,
  block/.style={draw=OilBorder, rounded corners=2pt, fill=PaperNoteBack, align=center, minimum height=0.78cm, text width=3.0cm, font=\small\color{OilDark}},
  driver/.style={draw=OilBorder, rounded corners=2pt, fill=white, align=center, minimum height=0.78cm, text width=3.1cm, font=\small\bfseries\color{OilDark}},
  damp/.style={draw=OilBorder, rounded corners=2pt, fill=white, align=center, minimum height=0.78cm, text width=3.1cm, font=\small\bfseries\color{OilDark}},
  arrow/.style={-{Latex[length=2.4mm]}, line width=0.55pt, draw=OilGray},
  plus/.style={font=\scriptsize\bfseries\color{OilDark}},
  minus/.style={font=\scriptsize\bfseries\color{OilDark}}
]
\node[driver] (shock) {供应中断\\剩余缺口 \(G_t\)};
\node[driver, right=of shock] (risk) {制度风险\\风险溢价};
\node[block, right=of risk] (price) {目标价格\\\(P_t^\ast\)};
\node[block, below=of price] (actual) {递推价格\\\(\hat P_t\)};
\node[damp, below=of shock] (buffer) {SPR/库存\\缓冲释放};
\node[damp, below=of risk] (route) {绕道运输\\需求收缩};
\node[damp, below=1.05cm of actual] (expect) {缓冲确认\\预期修复};

\draw[arrow] (shock) -- node[above, plus] {推高} (risk);
\draw[arrow] (risk) -- node[above, plus] {再定价} (price);
\draw[arrow] (price) -- node[right, plus] {调整} (actual);
\draw[arrow] (buffer) -- node[left, minus] {削峰} (shock);
\draw[arrow] (route) -- node[above, minus] {缩小缺口} (buffer);
\draw[arrow] (actual) -- node[right, plus] {高价反馈} (route);
\draw[arrow] (expect) -- node[right, minus] {降低溢价} (actual);
\draw[arrow] (buffer) -- node[below, minus] {可信度形成} (expect);
\draw[arrow] (route) -- node[left, minus] {修复预期} (expect);
\end{tikzpicture}
\caption{短期价格形成反馈回路}
\label{fig:tikz-feedback-loop}
\end{figure}

\subsubsection{物理供需层}

本文首先构造题面物理供需层。第 $t$ 日有效需求为

\begin{equation}
  \begin{aligned}
  \tilde D_t &= D_0\left(\frac{\hat P_{t-1}}{P_0}\right)^{\varepsilon_t},\\
  L_t &= \max(D_0-\tilde D_t,0),\\
  C_t &= \max(C_t^{\mathrm{obs}}-\rho L_t,0),\quad \rho=0.35,\\
  D_t &= \max(\tilde D_t-C_t,0.70D_0).
  \end{aligned}
\end{equation}

其中 \(\tilde D_t\) 为价格弹性调整后的需求，\(L_t\) 为价格弹性已解释的需求下降，\(C_t^{\mathrm{obs}}\) 为题面给出的需求下降量递推值。为避免价格弹性和题面需求下降重复全额扣减，本文仅将 \(35\%\) 的价格弹性下降视作与观测需求下降重叠，其余部分表示非价格行为调整和强制消费削减。第 $t$ 日不含库存的有效供给为

\begin{equation}
  S_t^{(0)}
  =
  S_0 - I + R_t + Q_t.
\end{equation}

商业库存缓冲量 $B_t$ 受剩余缺口、价格压力、库存日释放上限和剩余库存约束。记
\(\omega_t\in[0,1]\) 为由缺口压力和价格压力共同生成的库存释放压力，则

\begin{equation}
  B_t
  =
  \min\left\{
    \lambda(0.65+0.35\omega_t)\max(D_t-S_t^{(0)},0),\,
    \bar B_t(K_{t-1}),\,
    K_{t-1}
  \right\}.
\end{equation}

最终剩余缺口为

\begin{equation}
  G_t = \max(D_t-S_t^{(0)}-B_t,0).
\end{equation}

为避免将政策与物流缓冲写成纯时间函数，本文进一步定义当日客观缓冲覆盖率

\begin{equation}
  \eta_t^{\mathrm{obs}}
  =
  \frac{R_t+Q_t+B_t}
  {\max\{D_t-(S_0-I),\,1\}},
\end{equation}

其中分子表示 SPR、绕道运输和商业库存共同形成的当日缓冲供给，分母表示不考虑缓冲时的总短缺压力。考虑到库存、运输和政策执行信息并非当日完全透明，市场用于定价的覆盖率采用一阶感知更新：

\begin{equation}
  \eta_t
  =
  \eta_{t-1}
  +
  \alpha_\eta\left(\eta_t^{\mathrm{obs}}-\eta_{t-1}\right),
  \quad \alpha_\eta=0.35 .
\end{equation}

$\eta_t$ 越高，说明市场越相信缺口正在被覆盖，缓冲确认和预期修复越容易形成。实际计算中将覆盖率限制在合理区间内，避免极小分母造成异常放大。

\subsubsection{市场价格形成层}

原油期货价格不仅反映当日物理缺口，也反映市场对冲突持续时间、政策缓冲可信度和运输修复进度的预期。因此本文在物理供需层之上增加市场价格形成层。

为降低方程结构的任意性，本文采用一个极简的异质性交易者线性化框架。参考 Kilian（2009）、Hamilton（2009）以及 Kilian 与 Murphy（2014）关于物理供给、预防性需求、库存和投机性需求的分解，假设期货市场的边际定价者由两类主体构成：一类是基于炼厂采购、套期保值和现货短缺进行定价的供需交易者，另一类是基于封锁持续时间、航运保险、政策协调和风险偏好进行定价的风险交易者。令第 \(t\) 日相对目标价格偏离满足局部线性近似：

\begin{equation}
  \frac{P_t^\ast-P_0}{P_0}
  \approx
  a_G\frac{G_t/D_0}{|\varepsilon_t|}
  +a_B\frac{I}{D_0}Z_t
  +a_U U_t
  +a_F F_t
  -a_H\frac{H_t}{P_0}
  -a_L\frac{L_t}{P_0}.
\end{equation}

其中 \(G_t\) 表示剩余物理缺口，\(Z_t\) 表示封锁风险被市场吸收后的有效强度，\(U_t\) 表示持续封锁带来的制度不确定性，\(F_t\) 表示恐慌情绪强度，\(H_t\) 与 \(L_t\) 分别表示缓冲确认和预期修复带来的折价。这个表达式可理解为对两类交易者超额需求函数在战前均衡点附近做一阶展开：物理缺口和风险预防性需求推高价格，政策可信度和运输修复预期降低风险溢价。因此，后续目标价格的线性加总并非任意拼接，而是局部均衡价格的一阶近似。

具体时间函数来自三个建模约束。第一，剩余物理缺口和风险溢价分开建模，以对应油价冲击中文献强调的供给、预防性需求和库存渠道。第二，题面给出恐慌情绪放大持续 1--2 周，故风险溢价的建立项采用 \(1-e^{-t/7}\) 表示约一周尺度的信息吸收，指数衰减项表示市场在观察到 SPR、库存和绕道修复后逐步消化制度风险。第三，政策缓冲并非瞬时被市场完全相信，缓冲确认折价和预期修复折价分别刻画政策可信度形成与运输修复预期对目标价格的向下修正。新版短期冲击模型不再让两类折价只由日历时间决定，而是令其受到市场感知覆盖率 \(\eta_t\) 的平滑激活：当 SPR、库存和绕道对总短缺的覆盖率持续提高时，折价项才逐步生效。由此得到的指数项对应“冲击建立--风险消化--缓冲确认”三类时间过程。

\begin{equation}
\begin{aligned}
  P_t^\ast
  ={}& P_0
  + \phi P_0 \frac{G_t/D_0}{|\varepsilon_t|}
  + P_0\beta\frac{I}{D_0}
    \left(1-e^{-t/7}\right)e^{-\kappa_B t} \\
  & + P_0U\left(1-e^{-t/18}\right)
  + 0.45P_0F_t
  - H_t - L_t .
\end{aligned}
\end{equation}

其中 $\phi$ 为物理缺口向期货目标价格传导的价格形成系数，对应上式中的 \(a_G\)，不替代题面需求弹性；$\kappa_B=0.004$ 为封锁风险溢价的日度衰减速度，表示市场逐步吸收封锁制度风险的速度；$H_t=H(\eta_t)$ 为缓冲确认折价，$L_t=L(\eta_t,\Delta\eta_t)$ 为预期修复折价。也就是说，折价项不再是外生削减价格的自由项，而是由可观测的 SPR、库存和绕道覆盖率驱动。这里的 \(\varepsilon_t\) 仍承担题面低弹性约束，\(\phi\) 只刻画期货市场对剩余缺口的折价传导强度。模型递推式为

\begin{equation}
  \hat P_t
  =
  \hat P_{t-1}
  + v(P_t^\ast-\hat P_{t-1})
  + \gamma_F(F_t-F_{t-1}),\quad \gamma_F=2.5 .
\end{equation}

该结构保留了题面中的离散递推思想，同时允许期货市场对风险和预期做再定价。恐慌动量项 \(\gamma_F(F_t-F_{t-1})\) 只刻画恐慌强度快速变化时的一次性日度调整，不替代恐慌溢价本身；若恐慌强度稳定，则该项为零。风险溢价、缓冲确认折价和预期修复折价均需要同时满足经济含义清晰、方向可解释、消融后误差明显恶化以及扰动后结论保持稳定四个条件，才保留在主模型中。

从微观直觉看，式中的物理缺口项对应炼厂和贸易商面对短缺时愿意支付的边际溢价；风险项对应持有库存、安排绕道运输和承担航运保险风险所要求的补偿；折价项对应政策缓冲被观察到后，市场降低预防性持有需求和投机性风险溢价。本文没有直接建立完整的微观效用最大化模型，而是采用与能源冲击文献一致的结构化简形式，并通过基准对比、机制消融和扰动检验约束其自由度。

\begin{paperalgobox}{算法 1\quad 短期冲击模型的日度更新流程}
\textbf{输入：}初始价格 \(P_0\)，基准供需 \(D_0,S_0\)，供应中断 \(I\)，SPR、库存、绕道、恐慌与风险参数。\\
\textbf{输出：}冲突窗口内的日度模拟价格 \(\hat P_t\) 及机制贡献序列。

\begin{enumerate}[leftmargin=2.2em,itemsep=0.25em,topsep=0.45em,label=\arabic*.]
  \item 初始化 \(\hat P_0=P_0\)，商业库存 \(K_0\)，恐慌强度 \(F_0\)。
  \item 对 \(t=1,\ldots,T\) 逐日递推：
  \item 根据上一日模拟价格和短期弹性计算有效需求 \(D_t\)。
  \item 计算不含库存的有效供给 \(S_t^{(0)}=S_0-I+R_t+Q_t\)。
  \item 根据剩余缺口、价格压力、日释放上限和剩余库存计算商业库存缓冲 \(B_t\)。
  \item 得到剩余缺口 \(G_t=\max(D_t-S_t^{(0)}-B_t,0)\)。
  \item 计算客观缓冲覆盖率 \(\eta_t^{\mathrm{obs}}\)，再通过一阶感知更新得到 \(\eta_t\)，并由 \(\eta_t\) 激活缓冲确认和预期修复折价。
  \item 由 \(G_t\)、风险溢价、制度不确定性、恐慌强度、覆盖率驱动的缓冲确认和预期修复计算目标价格 \(P_t^\ast\)。
  \item 递推更新
  \[
    \hat P_t=\hat P_{t-1}+v(P_t^\ast-\hat P_{t-1})+\gamma_F(F_t-F_{t-1}) .
  \]
  \item 记录物理缺口、风险溢价、政策缓冲和预期折价的边际贡献。
\end{enumerate}
\end{paperalgobox}
